\documentclass[11pt,letterpaper]{article}

\usepackage[margin=1in]{geometry}
\usepackage{amsmath,amssymb}
\usepackage{booktabs}
\usepackage{setspace}
\usepackage{natbib}
\usepackage[hidelinks]{hyperref}
\usepackage{array}

\bibpunct{(}{)}{;}{a}{}{,}
\onehalfspacing
\setlength{\parskip}{0.4em}

%% ---------------------------------------------------------------------------
%% SSRN requires the title and ALL authors with affiliations to appear ON the
%% PDF. Fill these two lines in before uploading. (The Management Science
%% submission uses a SEPARATE, fully anonymised build -- do not upload this one
%% for double-anonymous review.)
%% ---------------------------------------------------------------------------
\newcommand{\AUTHORNAME}{[AUTHOR NAME]}
\newcommand{\AUTHORAFFIL}{[AFFILIATION]}
\newcommand{\AUTHOREMAIL}{[EMAIL]}

\title{\textbf{Alpha Without a Peer Group:\\
Self-Referential Factor Investing in Perpetual Futures}}
\author{\AUTHORNAME\\ \small \AUTHORAFFIL\\ \small \texttt{\AUTHOREMAIL}}
\date{This version: \today}

\begin{document}
\maketitle

\begin{abstract}
\noindent
Cross-sectional factor investing rests on an assumption so basic it is rarely
stated: that assets can meaningfully be ranked against one another. Equities
supply the anchor that makes this reasonable---shared accounting, industry
structure, and cash flows. Cryptocurrency perpetual futures supply none of it.
We ask what happens to factor construction when the comparability anchor is
absent. Transporting ten factors from the published equity literature to 112
perpetual contracts over a 363-week panel, we find that ranking each contract
against its own trailing history dominates ranking it against its peers in nine
of eleven cases, improving mean annualised Sharpe from 1.01 to 1.45. We
decompose the effect: approximately 41\% is attributable to scale normalisation,
the remainder to distributional robustness, since a rank transform is invariant
to distributional shape while standardisation assumes approximate normality. We
further show that preserving factor independence through portfolio
construction---forming one book per factor and combining returns rather than
blending scores---adds 0.665 Sharpe ($t = 5.82$) when both constructions are
tuned over an identical grid. We introduce a factor exploiting aggressor-side
information that centralised crypto venues publish exactly and equity research
must infer, which ranks third of eleven by leave-one-out contribution. Results survive walk-forward selection across seven rebalance
anchors (mean out-of-sample Sharpe 2.068, all profitable), a held-out universe
test with a matched-breadth control, transaction costs at four times the
modelled level, and a deflated Sharpe correction whose trial count is read from
an executed experiment log rather than recalled.

\medskip
\noindent\textbf{Keywords:} factor investing, cross-sectional ranking,
cryptocurrency, perpetual futures, portfolio construction, backtest overfitting

\noindent\textbf{JEL:} G11, G12, G14, C58

\medskip
\noindent\fbox{\parbox{0.96\textwidth}{\small
\textbf{[PLACEHOLDER --- COMPLETE BEFORE UPLOAD]} SSRN requires a disclosure
statement, displayed with the abstract, where AI tools were used in preparing
the work. Insert that statement here.}}
\end{abstract}

\thispagestyle{empty}
\newpage
\setcounter{page}{1}

%% ---------------------------------------------------------------------------
\section{Introduction}

Every cross-sectional factor model begins with a comparison. We say a stock is
cheap, or small, or heavily traded, and each of those words carries an implicit
completion: \emph{relative to other stocks}. The comparison feels natural
because equity markets supply the scaffolding that makes it meaningful. Two
firms report earnings under common accounting standards, operate in classifiable
industries, and are valued against cash flows that exist. When we rank them, we
rank like against like.

Cryptocurrency perpetual futures dissolve that scaffolding. There are no
earnings, no industry classification, and no accounting identity tying price to
anything. Asking whether one contract's turnover is high \emph{relative to
another contract's} is a question the data cannot really answer, because the two
quantities were never commensurate. Yet the factor literature, when it has
turned to this market, has largely imported the cross-sectional apparatus
unchanged.

This paper asks what that assumption costs, and what replaces it.

Our approach is deliberately conservative. Rather than proposing new signals, we
take ten factors established in the published equity literature and transport
them unchanged to 112 perpetual contracts across a 363-week panel. The factors
retain their original definitions and sign conventions; we fit nothing. The only
quantity we vary is the frame of comparison. In one specification each contract
is ranked against its peers in the same week, as the literature does. In the
other, each contract is ranked against its own trailing history---a
self-referential frame asking not ``is this asset unusual compared to others''
but ``is this asset unusual \emph{for itself}.''

The self-referential frame wins in nine of eleven cases, raising mean annualised
Sharpe from 1.01 to 1.45. It does not merely win; we can say why. A controlled
intermediate specification---standardising each asset against its own history
before ranking cross-sectionally---recovers 41\% of the gap, identifying scale
heterogeneity as one mechanism. The residual survives standardisation, and we
attribute it to distributional robustness: a rank transform is invariant to the
shape of the distribution it is applied to, whereas a $z$-score assumes
approximate normality, an assumption these distributions violate severely.

A second finding concerns what happens after the ranking. Standard practice
blends factor scores into a composite and forms a single portfolio from the
ordering it induces. The operation is lossy in a specific way: several
near-independent signals are compressed into one ranking, and their independence
is discarded before any position is taken. Forming one book per factor and
combining the \emph{returns}, holding every hyperparameter fixed, improves
performance by 0.665 Sharpe ($t = 5.82$), winning in 15 of 17 matched
comparisons. Both constructions were tuned over an identical grid; comparing a
tuned method against an untuned baseline is the most common way this literature
manufactures a result, and we take care not to.

Third, and connected directly to our thesis, we construct a factor that depends
on information this market publishes and equity markets do not. Research on
Chinese A-shares established that predictive content lies in the \emph{shape} of
the per-bar trade-size distribution rather than its mean. That measure is
necessarily direction-blind: A-share tick data does not reliably identify the
initiating side, so an unusually large purchase and an unusually large sale are
indistinguishable within the distribution. Perpetual futures publish taker-side
volume on every bar, permitting the same distribution to be partitioned by
initiating side and its asymmetry measured. We are careful not to overstate
this: US equity researchers can \emph{infer} direction via the tick rule
\citep{leeready1991}, with known error rates. The distinction is measured versus
estimated, not possible versus impossible---but it is real, and it illustrates
that what differs across these markets is not only the assets but the
information environment.

Finally, this paper is unusual in what it reports about itself. During
preparation, an audit of every headline claim against the code that supposedly
produced it identified four figures with no executable artifact, including a
pair of statistics that proved on inspection to be tail-index values transcribed
from a document concerning an unrelated strategy. We report this in Section~7,
together with the infrastructure now preventing recurrence: an append-only
experiment registry from which the multiple-testing correction reads its inputs,
and an automated gate asserting that the deployed system computes the same
positions as the research code. The relationship between a reported backtest and
the code that generated it is, we argue, an empirical question that ought to be
tested rather than assumed.

%% ---------------------------------------------------------------------------
\section{The comparability assumption}

Let $f_{i,t}$ denote a raw factor observation for asset $i$ at time $t$. A
cross-sectional model transforms this into a score by ranking within a date,
$s^{\mathrm{XS}}_{i,t} = \mathrm{rank}\!\left(f_{i,t} \mid \{f_{j,t}\}_{j=1}^{N}\right)$.
This is meaningful only if the $f_{j,t}$ are commensurate---if comparing asset
$i$'s value against asset $j$'s carries information about their relative future
returns. In equities this holds approximately, sustained by common exposure to
observable fundamentals \citep{famafrench1993} and by characteristics with
shared economic interpretation \citep{danieltitman1997}.

The self-referential alternative ranks within an asset instead,
$s^{\mathrm{TS}}_{i,t} = \mathrm{rank}\!\left(f_{i,t} \mid \{f_{i,\tau}\}_{\tau=t-w}^{t-1}\right)$,
making no cross-asset comparison at all. Its closest antecedent is time-series
momentum \citep{moskowitz2012}, which showed that an asset's own past return
predicts its future return independently of any cross-sectional ordering. Our
contribution is not the primitive but its systematic application across eleven
factors spanning price, flow, positioning and trade-size distribution, together
with identification of the mechanism.

The two transformations degrade differently. As the cross-section becomes more
heterogeneous, $s^{\mathrm{XS}}$ increasingly ranks \emph{asset identity} rather
than \emph{signal}: a contract whose order flow is systematically larger than
another's occupies the same tail of the cross-sectional distribution week after
week, whether or not anything informative has occurred. $s^{\mathrm{TS}}$ is
immune by construction, since its comparison set contains only the asset itself.

We are explicit about what the evidence supports. We find \emph{no sign
reversal}: cross-sectional ranking remains profitable, at mean annualised Sharpe
1.011 across eleven single-factor books. The claim is degradation, not failure.
A stronger reading---that the cross-sectional framework does not transfer to
markets lacking a comparability anchor---is not supported by these data, and we
do not make it.

%% ---------------------------------------------------------------------------
\section{Data and setting}

\textbf{Universe.} 112 USDT-margined perpetual futures contracts, screened for
continuous trading and finite positive prices, over a 363-week panel. After
factor warm-up the strategy is evaluable on 308 weekly rebalances and active on
282.

\textbf{Sources.} Price and volume are constructed from hourly bars drawn from a
public bulk archive and resampled weekly. Intraday inputs are reduced from
1-hour and 5-minute bars to one record per UTC day. Positioning data---open
interest and long/short ratios---comes from a keyless exchange metrics endpoint
at daily frequency. No proprietary or paid data is used.

\textbf{Why perpetual futures.} Two properties are load-bearing. First,
aggressor side is published: every bar carries taker-buy volume, where equity
microstructure research must infer direction \citep{leeready1991}. Second,
funding is observable and material: a price-only backtest of a long/short book
in this market overstates returns by approximately 3.9\% annually. We subtract
realised funding at every rebalance and tilt the book away from positions it
must pay to hold.

\textbf{A resolution constraint discovered by failure.} The trade-size
distribution factors cannot be computed hourly---a day divides into roughly
twelve bars per side, below the twenty-observation minimum the shape estimator
requires. Measured computability is 0\% at 1-hour and 100\% at 5-minute. The
remaining intraday factors, however, measured better at hourly resolution. The
pipeline therefore runs at two resolutions simultaneously, not as an
optimisation but as a constraint the data imposed.

%% ---------------------------------------------------------------------------
\section{Factors}

\subsection{Design discipline}

Every factor is specified \emph{a priori} from its source, including its sign.
Nothing is fitted. This determines what can be overfit: with no estimated
parameters anywhere in the model, the entire overfitting surface reduces to a
handful of discrete hyperparameters, and Section~6.2 holds all of them out.

\subsection{The ten transported factors}

Each reproduces published equity research: abnormal turnover (negative),
smart-money VWAP ratio, realised signed jump, order-flow imbalance,
price--volume correlation in level and in instability, three positioning factors
capturing crowding among large accounts, and the kurtosis of the log trade-size
distribution.

One sign convention is inverted relative to its source. The smart-money factor
is documented as a reversal signal in A-shares; in perpetual futures it is
momentum-consistent. We report this as an empirical finding rather than a fitted
choice, and record it as a limitation: a single inverted sign, chosen with
knowledge of the data, is a channel through which information can leak.

\subsection{Directional trade-size asymmetry}

Consider one intraday bar containing $n$ trades totalling $v$ units of volume.
Its \emph{ticket size} is $v/n$, the typical size of an individual trade within
it. Across a day these form a distribution, and the A-share literature
establishes that its \emph{shape}---dispersion, skewness, kurtosis---carries
predictive content that its mean does not: a more concentrated, more
right-skewed distribution precedes stronger returns.

The limitation is that a right-skewed distribution is produced equally by a
large buyer and a large seller. The measure detects that unusually sized orders
occurred, but not who placed them.

With taker-buy volume $b_n$ published per bar, the distribution can be
partitioned by initiating side:
\[
\mathcal{B} = \{n : b_n > \tfrac{1}{2}v_n\}, \qquad
\mathcal{S} = \{n : b_n \le \tfrac{1}{2}v_n\},
\]
\[
\mathrm{TSKD}_t \;=\; \Delta\Big[\mathrm{skew}\big(\log \mathrm{ticket}_{n\in\mathcal{B}}\big)
\;-\; \mathrm{skew}\big(\log \mathrm{ticket}_{n\in\mathcal{S}}\big)\Big].
\]

A positive value indicates that unusually large trades cluster on the buying
side. The \emph{change} carries the signal where the level does not: what
predicts is not which side is currently large but which side has \emph{just
become} large.

Measured by leave-one-out on a matched 282-week sample, TSKD ranks \textbf{third
of eleven} factors by marginal contribution ($+0.046$); standalone under the
deployed configuration it returns annualised $1.442$ ($t = 3.35$), and the
contribution is present in \emph{both} halves of the sample ($+0.030$ and
$+0.079$), indicating structure rather than a single-period artefact. Two honest
qualifications: a quintile sort on TSKD is not monotonic (rank correlation
between quintile index and forward return $+0.296$), the traded extremes carrying
the signal while interior quintiles are noisy; and TSKD is \emph{not} unusually
orthogonal, with mean absolute correlation $0.374$ against a factor-average of
$0.301$. It earns its place through predictive content rather than
diversification.

%% ---------------------------------------------------------------------------
\section{Construction}

\textbf{One book per factor.} Rather than forming a composite score
$\bar{s}_{i,t} = K^{-1}\sum_k s^{(k)}_{i,t}$ and trading the ordering it
induces, we form a dollar-neutral quintile book per factor---long the top 20\%,
short the bottom 20\%, weights $\pm 1/n$---and combine the resulting return
series.

\textbf{Inverse-volatility combination.} Book returns are combined by inverse
trailing volatility, $w^{(k)}_t \propto 1/\hat\sigma^{(k)}_{t-1}$, normalised to
sum to one, with the volatility estimate lagged so no contemporaneous
information enters the weights. Where the risk model is not yet estimable the
book is held flat rather than traded equally weighted.

\textbf{Carry is a factor, not an adjustment.} Each score is penalised by the
contract's cross-sectional funding rank. This is the one place a cross-sectional
comparison is appropriate---funding is denominated identically across contracts,
so ``which contract costs more to hold'' is a genuine peer comparison. The
component is far more consequential than ``tilt'' suggests: removing it costs
$+1.048$ (TSKD), $+0.820$ (TKU), $+0.762$ (OFI), $+0.559$ (AVOL) and $+0.540$
(Quad) annualised Sharpe respectively. Traded alone on the matched sample carry
returns $1.250$ ($t = 2.91$). Carry is a first-order signal in perpetual futures,
every single-factor figure in this paper includes it, and we say so rather than
presenting those numbers as the ranking effect in isolation.

\textbf{Turnover limitation as cost insurance.} A uniform partial adjustment
caps per-rebalance turnover. At the modelled maker fee this \emph{costs}
performance monotonically (Table~\ref{tab:cost}). What it purchases is
robustness to cost misestimation: the uncapped book is superior only while
realised costs remain near the assumption, deteriorating sharply beyond roughly
four times it, while the capped book remains viable to eight times. The deployed
configuration therefore uses a cap that is not Sharpe-maximising---a deliberate
divergence we state rather than leave to be discovered.

%% ---------------------------------------------------------------------------
\section{Results}

\subsection{The comparison frame}

Table~\ref{tab:frame} reports eleven single-factor books under both ranking
frames, holding data, universe, costs and construction fixed. Because carry
contributes 0.5--1.0 Sharpe to every book, we report the comparison both as the
strategy trades it and with carry removed entirely, so the ranking effect is
visible in isolation. As traded, mean annualised Sharpe rises from 1.011 to
1.451 (nine of eleven factors improve); with carry fully removed it rises from
0.383 to 0.648 (eight of eleven). \textbf{The ranking effect is independent of
carry}---carry amplifies it but does not create it. \emph{No} factor changes
sign under either specification.

To identify the mechanism we introduce an intermediate specification:
standardise each asset against its own trailing window, then rank
cross-sectionally. This recovers 41\% of the gap (1.011 $\rightarrow$ 1.192
$\rightarrow$ 1.451), identifying scale heterogeneity as one mechanism but not
the dominant one. We attribute the residual to distributional robustness:
standardisation normalises only the first two moments, whereas a rank against an
asset's own history is invariant to the entire distributional shape.

\subsection{Construction}

Both constructions were swept over an identical hyperparameter grid (1,286
executed configurations), so a tuned method is compared against a tuned
baseline. Across 17 matched cells the proposed construction returns 1.657
against the conventional 0.992, a mean difference of $+0.665$ with
$t = 5.82$, favourable in 15 of 17 cells. Best-of-grid across all six
construction variants places the proposed construction first on both maximum
(2.571) and mean (1.537), the conventional last (1.971 and 0.995).

\subsection{Full strategy}

The deployed configuration returns annualised Sharpe 2.161 ($t = 5.03$) over 282
weekly rebalances, at mean gross exposure 1.011 and mean turnover 0.369.

%% ---------------------------------------------------------------------------
\section{Validation}

\subsection{Walk-forward with rebalance-anchor variation}

Configuration is selected on a trailing 104-week window and evaluated on the
following 26 weeks, rolling forward. Because a weekly strategy also carries an
unmeasured exposure to which weekday it rebalances \citep{hoffstein2018}, the
panel is rebuilt from hourly bars on each of seven anchors and the procedure
repeated (Table~\ref{tab:anchor}).

Mean out-of-sample Sharpe is 2.068 with all seven anchors profitable and 182
out-of-sample weeks per anchor. We report the mean rather than the Monday
figure, and note explicitly that Monday---the schedule actually traded---is the
most favourable of the seven. Six of seven anchors clear the $t > 3.0$ threshold
proposed by \citet{harvey2016}; Friday, at 2.95, does not. Results are
additionally stable across fold geometry: seven train/test configurations from
78/13 to 156/26 all produce profitable out-of-sample series with $t$ from 3.70
to 5.55, and the geometry used for the headline is mid-range rather than
maximal.

\subsection{Held-out universe}

Time-series testing does not establish that hyperparameters generalise across
\emph{assets}. We therefore hold out symbols: configuration is selected on 78
contracts and evaluated on 34 that had no influence on the choice, across eight
random splits, with the cross-section rebuilt within each side so ranks,
quintile boundaries and the cost model are computed only among symbols present.

A naive comparison here is misleading, and we flag it because we initially made
it. A 34-contract book is mechanically inferior to a 78-contract one, since
information ratio scales with the square root of breadth \citep{grinold1989}.
We therefore add a matched-breadth control: the same selected configuration
evaluated on a random 34 of the \emph{training} contracts.

\begin{center}
\begin{tabular}{lcl}
\toprule
comparison & ratio & interpretation \\
\midrule
held-out / train (78) & 51\% & conflates two effects; not interpretable \\
control (seen, 34) / train (78) & 62\% & breadth, mechanical \\
\textbf{held-out / control} & \textbf{83\%} & \textbf{overfitting, at matched breadth} \\
\bottomrule
\end{tabular}
\end{center}

All eight splits are profitable out of universe. More informative than the ratio
is the stability of the selection itself: all eight chose identical funding
tilt, turnover treatment and quintile width, and five of eight the same ranking
window. Fitted noise would not converge this way. We note the limitation: only
one of eight held-out splits reaches $t > 3.0$, a consequence of reduced breadth
rather than of the effect's absence, and the held-out and control distributions
overlap substantially. The defensible claim is that retention is materially
above 50\%, not that it is precisely 83\%.

\subsection{Multiple testing with a measured trial count}

The deflated Sharpe ratio \citep{bailey2014} requires the number of trials and
their dispersion. Both describe the \emph{search} rather than the strategy and
cannot be recovered from the winning backtest, so in practice they are supplied
from recollection---precisely the quantity the correction exists to discipline.
We instead log every evaluated configuration to an append-only registry, written
by the code performing the evaluation, and read both inputs from it.

Under the strategy-search trial set ($N = 211$) the deflated Sharpe is 0.9885.
Under the most conservative reading, which additionally counts the 1,075
construction-ablation cells executed \emph{after} the strategy was specified in
order to measure the baseline fairly, $N = 1{,}286$ and the value is 0.9475,
below the conventional threshold. We report both. Its instability is itself the
argument against the conservative reading: the value fell from 0.9507 to 0.9475
purely because an ablation continued running. A statistic that penalises
conducting additional controlled experiments is not measuring selection bias.

\subsection{Transaction-cost robustness}

Costs are modelled as liquidity-scaled maker fees, from 1 basis point for the
most liquid quintile to 5 for the least, interpolated on cross-sectional
dollar-volume rank. Table~\ref{tab:cost} scales the entire curve; the deployed
configuration remains viable at eight times the modelled cost.

\subsection{Research--production equivalence}

The relationship between a published backtest and the system trading it is
normally asserted. We test it. Research and live paths call a single shared
implementation, and an automated gate asserts point-in-time integrity,
determinism, dollar-neutrality, non-emptiness and identical position directions,
exiting non-zero so it can block deployment. Measured: look-ahead leak
0.00e+00 on every factor; maximum net exposure 9.02e-17; 308 rebalances built,
active on 282; \textbf{zero direction mismatches across 4,480 comparisons}.

The point-in-time assertion recomputes a historical score with future
observations deleted and requires exact reproduction. This test exists because it
previously \emph{failed}: a ranking helper applied along the wrong axis ranked
each observation against the asset's entire history, future included, altering
historical scores for 98 of 112 contracts. Research and live code agreed
perfectly throughout, because both called the same incorrect function.

%% ---------------------------------------------------------------------------
\section{What did not survive}

\subsection{An audit of our own claims}

Before submission we searched the repository for the executable artifact behind
every headline figure. Four had none: a walk-forward Sharpe range that proved to
be tail-index statistics from a document about an unrelated strategy; a
timing-luck dispersion with no script or output; a pair of deflated Sharpe values
with no artifact; and a conventional-construction baseline that was approximately
the conventional construction's \emph{mean} compared against the proposed
construction's \emph{maximum}. A separate claim---that one factor reverses sign
between ranking frames---was tested directly and failed: zero of eleven reverse.

We report this because the corrected figures are in this paper and the discarded
ones are not, and because the episode motivates the infrastructure in
Sections~6.3 and~6.5.

\subsection{Negative results}

We conjectured that crypto factors are structurally more orthogonal than equity
factors. Measured pairwise correlation among Alpha101-style books is $+0.2173$
in this market against a reported 0.159 in equities: factors here are
\emph{more} correlated, not less. The orthogonality that makes the per-factor
construction valuable derives from data-source diversity---price, flow,
positioning, trade-size distribution---not from the asset class.

One candidate factor produced a Spearman information coefficient of $+0.0506$
($t = 6.10$) alongside a Pearson coefficient of only $+0.0172$ and a decile table
with no monotonic gradient. Returns here are severely fat-tailed; a factor can
order the typical asset correctly while being wrong about the few observations
that dominate realised profit.

%% ---------------------------------------------------------------------------
\section{Practical implications and limitations}

Reconsider the comparison frame before reconsidering the signal: the improvement
documented here required no new data and no new factor, only a change in what
each observation is ranked against. Any market with heterogeneous constituents
and no fundamental anchor is a candidate for the same treatment. Do not blend
scores; forming separate books and combining their returns preserves
independence that is expensive to acquire. Model funding explicitly. Choose the
cost-robust configuration rather than the Sharpe-maximising one. And test the
equivalence of research and production, which is normally assumed and which in
this project twice proved false, silently.

\textbf{Limitations.} Our evidence covers one venue's data over roughly seven
years. Out-of-sample testing is conducted in time and across the cross-section,
but not across venues. Capacity is not estimated in currency terms; the maker-fee
assumption will not survive at institutional scale in less liquid contracts, and
Section~6.4 should be read as bounding rather than resolving that concern. One
factor's sign convention was inverted with knowledge of the data. The held-out
universe result rests on eight splits and carries substantial uncertainty.

%% ---------------------------------------------------------------------------
\section{Conclusion}

Factor investing imported into cryptocurrency perpetual futures inherits an
assumption this market does not support: that its constituents are comparable.
We find the assumption costly but not fatal---cross-sectional ranking remains
profitable, while ranking each contract against its own history performs
materially better in nine of eleven cases, for reasons we can partly decompose.
We further find that how signals are combined after ranking matters more than
any individual signal, and we exhibit a factor whose construction depends on
information this market publishes exactly and equity research must infer.

We would emphasise the structure of the evidence as much as any single result:
out-of-sample in time and across assets, replicated over seven rebalance
schedules and seven fold geometries, corrected for multiple testing using a
trial count read from an executed log, stress-tested to eight times modelled
costs, and verified to be computed identically by the deployed system. Each test
was capable of overturning the result. We also report the four claims that did
not survive our own audit, because the credibility of a backtest is inseparable
from the reproducibility of the pipeline that produced it.

%% ---------------------------------------------------------------------------
\newpage
\section*{Tables}
\addcontentsline{toc}{section}{Tables}

\begin{table}[h]
\centering
\caption{Ranking frame, eleven single-factor books. Annualised Sharpe, net of
funding and liquidity-scaled maker costs.}
\label{tab:frame}
\begin{tabular}{lccc}
\toprule
specification & cross-sectional & self-referential & gap \\
\midrule
as traded (with carry)   & 1.011 & \textbf{1.451} & $+0.440$ \\
\quad factors improved   & ---   & 9 of 11        &          \\
carry removed            & 0.383 & \textbf{0.648} & $+0.265$ \\
\quad factors improved   & ---   & 8 of 11        &          \\
sign reversals           & ---   & \textbf{0 of 11} &        \\
per-asset standardised   & 1.192 & & (41\% of gap recovered) \\
\bottomrule
\end{tabular}
\end{table}

\begin{table}[h]
\centering
\caption{Walk-forward out-of-sample Sharpe by rebalance anchor. Configuration
selected on trailing data only; 182 out-of-sample weeks per anchor.}
\label{tab:anchor}
\begin{tabular}{lccccccc|c}
\toprule
anchor & Mon & Sat & Sun & Tue & Wed & Thu & Fri & mean \\
\midrule
OOS Sharpe & 2.713 & 2.496 & 2.423 & 1.855 & 1.753 & 1.656 & 1.579 & \textbf{2.068} \\
$t$        & 5.08 & 4.67 & 4.53 & 3.47 & 3.28 & 3.10 & 2.95 & \\
\bottomrule
\end{tabular}
\end{table}

\begin{table}[h]
\centering
\caption{Transaction-cost robustness. The entire liquidity-scaled cost curve is
multiplied by the stated factor.}
\label{tab:cost}
\begin{tabular}{lcccccc}
\toprule
cost multiple & 1$\times$ & 2$\times$ & 4$\times$ & 8$\times$ & 16$\times$ & 32$\times$ \\
\midrule
uncapped            & 2.364 & 2.247 & 2.014 & 1.547 & 0.618 & $-1.222$ \\
deployed (capped)   & 2.161 & 2.117 & 2.028 & 1.851 & 1.496 & 0.789 \\
\bottomrule
\end{tabular}
\end{table}

%% ---------------------------------------------------------------------------
\newpage
\begin{thebibliography}{99}

\bibitem[Asness et al.(2013)]{asness2013} Asness CS, Moskowitz TJ, Pedersen LH (2013) Value and momentum everywhere. \emph{Journal of Finance} 68(3):929--985.

\bibitem[Bailey and L\'opez de Prado(2014)]{bailey2014} Bailey DH, L\'opez de Prado M (2014) The deflated Sharpe ratio: Correcting for selection bias, backtest overfitting, and non-normality. \emph{Journal of Portfolio Management} 40(5):94--107.

\bibitem[Black and Litterman(1992)]{black1992} Black F, Litterman R (1992) Global portfolio optimization. \emph{Financial Analysts Journal} 48(5):28--43.

\bibitem[Daniel and Titman(1997)]{danieltitman1997} Daniel K, Titman S (1997) Evidence on the characteristics of cross sectional variation in stock returns. \emph{Journal of Finance} 52(1):1--33.

\bibitem[Fama and French(1993)]{famafrench1993} Fama EF, French KR (1993) Common risk factors in the returns on stocks and bonds. \emph{Journal of Financial Economics} 33(1):3--56.

\bibitem[Grinold(1989)]{grinold1989} Grinold RC (1989) The fundamental law of active management. \emph{Journal of Portfolio Management} 15(3):30--37.

\bibitem[Gu et al.(2020)]{gu2020} Gu S, Kelly B, Xiu D (2020) Empirical asset pricing via machine learning. \emph{Review of Financial Studies} 33(5):2223--2273.

\bibitem[Harvey et al.(2016)]{harvey2016} Harvey CR, Liu Y, Zhu H (2016) \ldots and the cross-section of expected returns. \emph{Review of Financial Studies} 29(1):5--68.

\bibitem[Hoffstein et al.(2018)]{hoffstein2018} Hoffstein C, Sibears J, Faber N (2018) Rebalance timing luck: The difference between hired and fired. \emph{Journal of Index Investing}.

\bibitem[Kyle(1985)]{kyle1985} Kyle AS (1985) Continuous auctions and insider trading. \emph{Econometrica} 53(6):1315--1335.

\bibitem[Lee and Ready(1991)]{leeready1991} Lee CMC, Ready MJ (1991) Inferring trade direction from intraday data. \emph{Journal of Finance} 46(2):733--746.

\bibitem[Ledoit and Wolf(2004)]{ledoit2004} Ledoit O, Wolf M (2004) A well-conditioned estimator for large-dimensional covariance matrices. \emph{Journal of Multivariate Analysis} 88(2):365--411.

\bibitem[Moskowitz et al.(2012)]{moskowitz2012} Moskowitz TJ, Ooi YH, Pedersen LH (2012) Time series momentum. \emph{Journal of Financial Economics} 104(2):228--250.

\end{thebibliography}

\end{document}
